\label{sec:identification}

The central identification challenge for all four O-series papers is establishing
causal rather than correlational relationships between map geometry and democratic
outcomes. This section describes the identification strategies used across the series,
with particular emphasis on the features that make redistricting cycles well-suited
as quasi-experiments.

\subsection{Redistricting Cycles as Natural Experiments}

The decennial redistricting cycle generates a near-ideal natural experiment for
studying the effects of map geometry on democratic outcomes. The experiment has
three favorable properties:

\textbf{Sharp discontinuity in treatment.} District boundaries change abruptly
following redistricting, creating a sharp discontinuity in treatment assignment
(which district a precinct belongs to) at a known calendar date. Precincts whose
boundaries are unchanged across the redistricting cycle serve as a natural comparison
group, while precincts that change district assignment receive the treatment.
This design is related to geographic regression discontinuity \citep{keeleTimmons2016}
but exploits temporal rather than spatial variation.

\textbf{Pre-redistricting baseline.} Outcomes from the two election cycles immediately
before redistricting (typically four House elections spanning eight years) provide
a pre-period baseline against which post-redistricting outcomes can be compared.
Because redistricting occurs at the start of each decade, a full two-term congressional
cycle is available in both the pre- and post-period under the same map.

\textbf{Plausible exogeneity of assignment mechanism.} While the decision to gerrymander
is endogenous, the specific boundary placement within a gerrymandered map is driven by
complex political calculations that include geographic constraints, VRA requirements,
incumbent protection, and intra-party negotiation. This complexity makes the specific
precinct-level district assignment difficult to fully anticipate from precinct-level
outcome data, supporting an approximate exogeneity assumption for precinct-level analysis
(distinct from the plan-level analysis, where endogeneity of gerrymandering is severe).

\subsection{Difference-in-Differences Design}

The primary identification strategy for O.1 and O.2 is a difference-in-differences (DiD)
design that compares changes in outcomes for precincts that experienced a large district
geometry change to precincts that experienced a small change, before and after redistricting.

Let $Y_{pt}$ be the outcome (turnout, competitiveness indicator) for precinct $p$ in
election year $t$. Let $\Delta\text{PP}_{p}$ be the change in the Polsby-Popper score
of the district containing precinct $p$ between the pre- and post-redistricting cycles.
The DiD estimator is:

\begin{equation}
Y_{pt} = \alpha + \beta \cdot \text{Post}_t + \gamma \cdot \Delta\text{PP}_p
  + \delta \cdot (\text{Post}_t \times \Delta\text{PP}_p) + \mathbf{X}_{pt}'\theta + \varepsilon_{pt}
\end{equation}

where $\text{Post}_t = 1$ for election years after redistricting, $\mathbf{X}_{pt}$ is
a vector of precinct-level covariates (median income, racial composition, urban-rural
classification), and $\delta$ is the parameter of interest: the effect of a one-unit
increase in Polsby-Popper change on the outcome change following redistricting.

Standard errors are clustered at the district level to account for within-district
correlation in outcomes, and at the state-year level to account for state-specific
electoral shocks.

\subsection{Instrumental Variables: The Algorithmic Instrument}

For outcomes where the DiD design is subject to selection concerns (e.g., if states
that gerrymander more aggressively also have systematically different political cultures
that affect turnout independently of redistricting), we supplement the DiD with an
instrumental variables design.

The instrument is the \textit{algorithmic counterfactual geometry}: the Polsby-Popper
score of the \textsc{bisect} district that contains each precinct in the algorithmic
plan, net of the enacted district Polsby-Popper. This instrument satisfies the
relevance condition by construction (algorithmic geometry is highly correlated with
actual geometry in states without gerrymandering, and systematically higher in
gerrymandered states) and the exclusion restriction by design (the algorithmic plan
is computed without reference to any outcome variable; partisan and turnout data
play no role in \textsc{bisect}'s partitioning).

The IV estimand recovers a local average treatment effect (LATE) for the subpopulation
of precincts whose district assignment is affected by the compactness difference between
the enacted and algorithmic plans---in practice, precincts near the district boundaries
where gerrymandering produces the most non-compact geometry.

\subsection{Event Study Design for Polarization (O.3)}

For legislative polarization, the DiD design must be applied at the legislator level
rather than the precinct level. The design uses redistricting-cycle transitions as
event dates: for each member of Congress, we record their DW-NOMINATE score in each
term and interact this with whether they entered the House after a redistricting cycle
that changed their district's geometry substantially.

The event study specification:

\begin{equation}
\text{DW-NOMINATE}_{it} = \alpha_i + \alpha_t
  + \sum_{\tau=-4}^{4} \beta_\tau \cdot \mathbb{1}[\text{cycle}_t = \tau] \cdot \Delta\text{PP}_i
  + \varepsilon_{it}
\end{equation}

where $\alpha_i$ is a member fixed effect, $\alpha_t$ is a Congress fixed effect,
$\tau$ indexes election cycles relative to the redistricting event, and $\Delta\text{PP}_i$
is the Polsby-Popper change for member $i$'s district in the redistricting event.
Coefficients $\beta_\tau$ for $\tau < 0$ serve as pre-trend tests; $\beta_\tau$ for
$\tau \geq 0$ estimate the dynamic treatment effect of compactness changes on polarization.
